\section{A New Scaling Dimension}

The next generation of AI systems will not be defined only by the capability of
their strongest model at a moment in time. They will also be defined by their
ability to sustain, adapt, and improve that capability over the life of a task.

Fumo Lab is building \systemname for this transition. Fusion is a
\emph{compound intelligence system}: it combines a changing portfolio of
models with an intelligence layer that learns how to understand the current
state, acquire the evidence required next, and allocate capability over time.
The ambition is not simply to compose more models. Composition describes a
system's structure; compounding describes how its capability grows.

\subsection{From model capability to trajectory capability}

On a bounded request, capability can often be approximated by the quality of one
response. In an agentic workflow, capability has to persist across a trajectory.
Every observation and action changes the state of the work. The system may need
to explore broadly, construct precisely, inspect the environment, challenge a
candidate, recover from failure, and verify completion---without losing the
intent and evidence accumulated along the way.

This changes what it means for an AI system to be capable. A strong answer at
one turn is not enough if later steps repeat work, act from incompatible
assumptions, spend compute on resolved questions, or declare success before the
result is verified. System capability includes the policy that decides what
intelligence the evolving task should receive at every transition.

In a bounded request, that policy can resemble a routing optimization. Across a
stateful trajectory, it becomes a persistent system capability. Better
allocation improves not only which model answers, but how intelligence is
assembled, preserved, and redirected as the task changes.

\subsection{Two coupled loops of progress}

Fusion advances through two connected loops. The first operates \emph{within}
a trajectory. The current task state exposes a decision-relevant uncertainty;
Fusion allocates a model--role or tool computation to acquire evidence; the
result updates the state; and the system allocates again. Allocation
intelligence governs this full loop rather than only its first request.

The second operates \emph{across} trajectories. Each completed task reveals how
Fusion interpreted the state, which computations it selected, what evidence
they produced, which evidence changed the commitment, and whether the outcome
succeeded. Those observations improve future task understanding, capability
estimation, allocation, and stopping.

Two forms of progress enter these loops:

\begin{equation}
  \text{system progress}
  \quad\longleftarrow\quad
  \underbrace{\text{model progress}}_{\text{supply}}
  \quad\times\quad
  \underbrace{\text{allocation progress}}_{\text{demand}}.
\end{equation}

This is an architectural claim, not a proposed empirical power law. Model
progress expands the frontier of what Fusion can acquire; allocation progress
improves how effectively that supply becomes an outcome.

\emph{Supply-side intelligence improvement} expands what Fusion can draw upon
as models advance in reasoning, specialization, reliability, speed, and cost.
\emph{Demand-side allocation intelligence} improves how effectively Fusion
turns that portfolio into useful evidence and action at each state. Either can
advance without waiting for the other; together, their effects can compound.

\begin{figure}[H]
\centering
\resizebox{0.88\textwidth}{!}{%
\begin{tikzpicture}[
  node distance=1.05cm and 1.25cm,
  box/.style={draw=fusionblue!70, rounded corners=3pt, very thick,
    fill=fusionblue!5, minimum height=0.92cm, minimum width=3.0cm,
    align=center, inner xsep=10pt},
  source/.style={draw=fusiongray!70, rounded corners=3pt, fill=black!2,
    minimum height=0.92cm, minimum width=3.0cm, align=center, inner xsep=10pt},
  arrow/.style={-{Latex[length=2.2mm]}, thick, draw=fusiongray}
]
\node[box] (state) {Current\\trajectory state};
\node[box, right=of state] (allocate) {Allocation\\intelligence};
\node[box, right=of allocate] (work) {Evidence / action};
\node[source, above=of state] (progress) {Model progress};
\node[source, above=of allocate] (supply) {Capability supply};
\node[source, below=of work] (outcome) {Trajectory outcome};
\node[source, left=of outcome] (learn) {Learning from behavior\\and feedback};

\draw[arrow] (progress) -- (supply);
\draw[arrow] (supply) -- (allocate);
\draw[arrow] (state) -- (allocate);
\draw[arrow] (allocate) -- (work);
\draw[arrow] (work.north) -- ++(0,0.42) -| (state.north);
\node[font=\small, fill=white, inner sep=1.5pt, above=0.27cm of allocate]
  {updated state};
\draw[arrow] (work) -- (outcome);
\draw[arrow] (outcome) -- (learn);
\draw[arrow] (learn) -- (allocate);
\end{tikzpicture}
}
\caption{Model progress expands the intelligence available to Fusion.
Allocation progress improves the inner trajectory loop and learns from outcomes
across trajectories.}
\label{fig:compounding-loops}
\end{figure}

The outer loop is recursive in a bounded, system-level sense. Fusion uses traces
of its own behavior and outcomes to improve how it assembles intelligence on
the next task. This does not imply unrestricted self-modification: policies
remain versioned, evaluated, constrained, and reversible. Section~\ref{sec:learning}
describes the learning mechanism and its safeguards in detail.

\subsection{A compounding performance frontier}

The external manifestation of this scaling dimension has two parts. On bounded
problems, Fusion should move the performance--compute frontier: comparable
quality with less frontier computation, and greater quality at comparable
compute. Across agentic trajectories, it should preserve effectiveness as the
state changes: higher end-to-end completion, fewer inconsistent transitions,
and lower compute per successful task.

This is not a vision of ``cheap AI.'' Cost reveals whether the system
understands its own work. Redundant frontier calls indicate that it cannot
distinguish available evidence from missing evidence. Premature use of a weaker
source indicates that it cannot recognize where stronger judgment is
load-bearing. Allocation intelligence must improve both capability and the
productivity of compute.

\principle{Model progress expands the supply of intelligence. Fusion compounds
it through allocation intelligence---within a task and across tasks.}
